% NDOT 论文详尽汇总
\documentclass[journal,12pt,onecolumn]{IEEEtran}
\IEEEoverridecommandlockouts
\usepackage{cite}

\usepackage{amsmath,amssymb,amsfonts}
\usepackage{algorithmic}
\usepackage{algorithm}
%\usepackage{algpseudocode}
\usepackage{graphicx}

\usepackage{booktabs}

\usepackage{textcomp}
%\usepackage[marginal]{footmisc}
\usepackage{xcolor}
\usepackage{subfigure}
\graphicspath{{./img/}}
\usepackage{tabularx}
\usepackage{makecell}
\usepackage{caption2}
%\usepackage{fontspec}
%\setmainfont{TeX Gyre Termes}
\usepackage{color}

\usepackage[UTF8]{ctex}

\begin{document}
\title{NDOT: Neuronal Dynamics-based Online Training for Spiking Neural Networks \\ 论文详细汇总}
\author{}
\date{}
\maketitle

\section{摘要}
本文基于原文 \texttt{`Algorithm/NDOT/9417\_NDOT\_Neuronal\_Dynamics\_ba.pdf`} 提取并整理，详细介绍了 NDOT（Neuronal Dynamics-based Online Training）方法的核心思想、关键公式、算法流程和优势。NDOT 通过神经元动力学表示连续时间依赖，实现前向时间在线训练，并避免传统 BPTT 的高内存开销。本文还总结了 NDOT 在实验中的表现和结论。\cite{jiang2024ndot}

\section{研究背景}
脉冲神经网络（SNN）通过离散脉冲传播信息，具备能效与时序敏感优势，但其二值脉冲生成函数不可导，使得直接训练面临困难。目前典型方法包括：
\begin{itemize}[leftmargin=*]
  \item ANN-to-SNN 转换：先训练 ANN 再映射到 SNN，但通常需要更多时间步来维持精度；
  \item BPTT + surrogate gradient：将 SNN 视为二值 RNN，沿时间展开计算梯度，内存消耗随时间步线性增长；
  \item 脉冲表示方法：用平均发放率等表征近似时间依赖，降低时序计算但可能增加延迟。
\end{itemize}
在此背景下，NDOT 提出了一种面向在线学习的前向时间训练框架，直接从神经元动力学中提取连续时间依赖。

\section{NDOT 方法}
\subsection{模型与梯度分解}
NDOT 的核心思想是将全梯度分解为空间梯度和时间梯度两部分。在经典 BPTT 中，权重梯度具有如下时间展开形式：
\begin{equation}
  \frac{\partial L}{\partial W^l} = \sum_{t=1}^T \frac{\partial L}{\partial s^l[t]} \frac{\partial s^l[t]}{\partial u^l[t]} \left(\frac{\partial u^l[t]}{\partial W^l} + \sum_{k<t} \prod_{i=k}^{t-1} \varepsilon^l[i] \frac{\partial u^l[k]}{\partial W^l}\right),
\end{equation}
其中
\begin{equation}
  \varepsilon^l[t] = \frac{\partial u^l[t+1]}{\partial u^l[t]} + \frac{\partial u^l[t+1]}{\partial s^l[t]} \frac{\partial s^l[t]}{\partial u^l[t]}
\end{equation}
表示离散时间依赖。

NDOT 采用神经元动力学的连续时间依赖表示：
\begin{equation}
  e^l[t] = \frac{u^l[t]-V_{th}s^l[t]}{u^l[t-1]-V_{th}s^l[t-1]}.
\end{equation}
这里 $u^l[t]$ 为第 $l$ 层第 $t$ 个时间步的膜电位，$s^l[t]\in\{0,1\}$ 为脉冲输出。

基于该表示，NDOT 将全梯度写成：
\begin{equation}
  \nabla_{W^l}L = \sum_{t=1}^T g^u_l[t] \, \hat{a}^{l-1}[t]^\top,
\end{equation}
其中空间梯度定义为：
\begin{equation}
  g^u_l[t] = \left(\frac{\partial L}{\partial s^l[t]} \frac{\partial s^l[t]}{\partial u^l[t]}\right)^\top.
\end{equation}
时间梯度 $\hat{a}^{l-1}[t]$ 则表示前一层脉冲历史的累积影响。

\subsection{时间梯度的前向跟踪}
NDOT 使用以下方式跟踪时间梯度：
\begin{align}
  \hat a^{l-1}[t] &= s^{l-1}[t] + \sum_{k<t} P^{l-1}_{k,t} \odot s^{l-1}[k], \\
  P^{l}_{k,t} &= \prod_{i=k}^{t-1} e^l[i] = \frac{u^l[t-1]-V_{th}s^l[t-1]}{u^l[k-1]-V_{th}s^l[k-1]}.
\end{align}
其中 $\odot$ 表示逐元素乘积。

该累积表达式可以递归地写成：
\begin{equation}
  \hat a^{l-1}[t] = e^{l-1}[t-1] \odot \hat a^{l-1}[t-1] + s^{l-1}[t].
\end{equation}
这一前向递推关系使得 NDOT 能在时间维上逐步更新，而无需回溯整个历史。

\subsection{NDOT 算法流程}
\begin{algorithm}[htbp]
\caption{NDOT 一次迭代}
\begin{algorithmic}[1]
  \STATE Input: 训练样本 $(x,y)$, 时间步数 $T$。
  \STATE Initialize $\hat a^l[0]=0$ 和 $e^l[0]=0$。
  \FOR{$t = 1, \dots, T$}
    \FOR{$l = 1, \dots, N$}
      \STATE 前向计算膜电位 $u^l[t]$ 和脉冲 $s^l[t]$。
      \STATE 计算连续时间依赖 $e^l[t] = \frac{u^l[t]-V_{th}s^l[t]}{u^l[t-1]-V_{th}s^l[t-1]}$。
      \STATE 递推时间梯度 $\hat a^l[t] = e^l[t-1] \odot \hat a^l[t-1] + s^l[t]$。
    \ENDFOR
    \FOR{$l = N, \dots, 1$}
      \STATE 计算空间梯度 $g^u_l[t] = (\partial L / \partial s^l[t])(\partial s^l[t] / \partial u^l[t])$。
      \STATE 计算即时权重梯度 $\nabla_{W^l}L[t] = g^u_l[t](\hat a^{l-1}[t])^\top$。
      \STATE 使用梯度优化器更新权重 $W^l$。
    \ENDFOR
  \ENDFOR
\end{algorithmic}
\end{algorithm}

\subsection{在线损失设计}
NDOT 采用即时损失 $L[t]$，使每个时间步的梯度可以独立计算，从而支持在线训练。论文强调这与基于平均发放率的离线损失不同，后者需要访问整个时间序列。

\section{NDOT 的优势}
\begin{itemize}[leftmargin=*]
  \item \textbf{前向时间在线学习}：梯度计算仅依赖当前时间步和已维护的前一层时间梯度，无需时间维上反向传播；
  \item \textbf{常数训练内存}：NDOT 的状态维护在时间维上为常数，与 BPTT 的 $O(T)$ 内存开销形成对比；
  \item \textbf{连续时间依赖建模}：通过 $e^l[t]$ 表示神经元的连续动态依赖，保持膜电位和放电后的历史影响；
  \item \textbf{兼容自动微分}：空间梯度 $g^u_l[t]$ 可通过现有深度学习框架自动求导计算；
  \item \textbf{直观解释}：论文通过 FTRL 视角说明 NDOT 以连续时间依赖隐式捕捉历史信息，弥补仅依赖即时损失方法的不足。
\end{itemize}

\section{实验验证}
论文在 CIFAR-10、CIFAR-100 和 CIFAR10-DVS 上验证了 NDOT 的性能，表明该方法在小时间步设置下仍能保持良好精度，并显著降低训练内存需求。同时，论文通过 GPU 内存占用对比，展示了 NDOT 的在线训练可行性。

\section{结论}
NDOT 通过神经元动力学的连续时间依赖和空间/时间梯度解耦，实现了可在线执行的前向时间 SNN 训练。该方法兼顾了梯度准确性和内存效率，是面向 neuromorphic 在线训练的有效方案。

\bibliographystyle{IEEEtran}
\bibliography{../SNNInformationReconstruction}

\end{document}
